\documentclass[a4paper,10pt]{article}
\usepackage[francais]{babel}
\usepackage{url}

\selectlanguage{francais}
\input{letter-ensiie.sty}

\selectlanguage{francais}
\enCouleur

\begin{document}


\name{St\'ephane \textsc{Genaud}}
\fonction{Directeur de l'antenne de Strasbourg de l'ENSIIE\\Professeur des Universités, ENSIIE}
\telephoneillkirch{+33 (0)3688 54542}
\email{genaud@ensiie.fr}

\begin{letter}[Avis concernant l'admission en L3 MIASHS de Florian BUCHHEIT]{à}%
{Jury d'Admission en Licence MIASHS}


Madame, Monsieur,

Je  connais Florian  BUCHHEIT depuis  son  intégration à  l'ENSIIE en  septembre
2014. L'ENSIIE est  une école publique d'Ingénieur en  Informatique sous tutelle
du  Ministère de  l'Enseignement  Supérieur et  de la  Recherche,  dont le  site
principal est à Evry.

Issu  des  classes  préparatoires  aux   grandes  écoles,  Florian  a  découvert
l'informatique à  l'ENSIIE. Pour l'avoir  eu également comme étudiants  dans des
cours de  programmation système et d'algorithmique,  je sais à quel  point cette
discipline  était nouvelle  pour Florian.   Le rythme  accéléré de  formation en
informatique est  difficile pour de  nombreux étudiants, dont  les connaissances
initiales sont inégales. Florian n'a que partiellement validé sa première année,
essentiellement pour des échecs sur les matières informatiques. Il a cumulé ce
retard sur sa deuxième année, ce qui a contraint l'école à prononcer un arrêt de
la scolarité car les 6 semestres d'enseignement ne peuvent être faits que dans
une période contrainte de temps de 8 semestres.

Cependant, je peux témoigner que Florian a toujours été un élève-ingénieur
extrêmement sérieux et assidu. J'ajoute que Florian est doté d'une grande
curiosité, l'amenant à démarrer des projets dans des domaines annexes pour y
appliquer ses connaissances en informatique (par exemple informatique musicale).

Je suis convaincu que Florian pourrait être un bon élément en licence MIASHS.
Son niveau en mathématiques est très bon, et d'autre part, je crois qu'il a mûri
son projet professionnel pendant ses deux années d'école d'ingénieur.

Veuillez agréer, Madame, Monsieur, l'expression de mes sentiments respectueux.
\end{letter}
\begin{flushright}
\includegraphics[width=.2\textwidth]{signgenaud.jpg}
\end{flushright}
\end{document}
